\section{Repairs, and a supplementary result}
\label{app:proofs}

\subsection{Statements we repaired}

While formalizing, we found several statements in our own earlier formulations
to be false. We replaced each with a true statement and give both here, so that
no reader mistakes the original for a result.
\begin{itemize}
\item \emph{Licences for dropping route memory.} The original said that routes
with the same boundary have equal endpoints \emph{only if} a normalization,
confluence or flatness licence holds. That is false: on a domain of three
routes, two have equal endpoints and no licence holds. We prove only
sufficiency (Theorem~\ref{thm:licences}).
\item \emph{Direct succession and data.} The original said that direct
succession excludes every data path. That is false: a genuine direct successor
can coexist with an independent data path. We prove the four exclusions and the
insufficiency of data reachability (Proposition~\ref{prop:direct-exclusions}).
\item \emph{Robustness.} The original said that every perturbation within the
model's margin preserves every advance and every required condition. That is
false. A correction by a quarter is not an integer difference. A receiver
that outputs a constant keeps every number positive but fails the checks. We
prove robustness in a rational extension, for one event and one box
(Proposition~\ref{prop:box}).
\item \emph{Unbounded maintenance.} The original said that a policy on a
finite invariant set yields an infinite run in a finite full state. That is
impossible, since no infinite run that strictly progresses fits a finite full
state. We prove it on a finite factor with unbounded registers
(Theorem~\ref{thm:unbounded}).
\item \emph{Reconstruction versus twin.} The original compared ancestry from
one target to another and is refuted for every prefix. We use the rooted form
(Proposition~\ref{prop:reconstruction}).
\item \emph{Identity.} The original made a claim about every identity
relation, and the relation that identifies everything refutes it. We prove
underdetermination for a stated class of identity relations
(Theorem~\ref{thm:identity}).
\end{itemize}

\subsection{Conditioning does not factor}

A natural shortcut for the joint behaviour of two branches is to condition on
a shared event and multiply the conditioned marginals. That is wrong in
general, and the fusion model's predecessor gives an exact case.

\begin{proposition}[Conditioning does not factor]
% claim: JOINT-COND
\label{prop:conditioning}
On the predecessor of the fusion model, a shared event has two outcomes, of
weights \snapGTwoSixConditioningEventWeights{} (ratio
\snapGTwoSixConditioningRelativeWeights{}). Conditioned on the event, and with the empty action word, the joint initial
observation associated with the lower-weight outcome has probability
\snapGTwoSixConditionedJointMass{}. Its two components are the corresponding
initial observations of the two branches. Each has conditioned probability
\snapGTwoSixConditionedLeftMass{}, and their product is
\snapGTwoSixProductOfConditionedMarginals{}. The correct joint
law is a weighted sum of products of the branch laws over fixed worlds.
\end{proposition}
